\newif\ifglobal

%%%%%%%%%%%%%%%%%%%%%%%
%%% Compile to view %%%
%%%%%%%%%%%%%%%%%%%%%%%

%%%%%%%%%%%%%%%%%%%%%%%%%%%%%%%%%%%%%%%%%%%%%%%%%%%
%%% Readme:                                     %%%
%%% https://www.overleaf.com/read/bwdjvfjksvjy/ %%%
%%%%%%%%%%%%%%%%%%%%%%%%%%%%%%%%%%%%%%%%%%%%%%%%%%%

% >>> M?: marks a module setting number or important notice

% >>> M4: Set {./relative-path} to external files for this module <<<
\def \modulefiles {./06-VDMA}
% To add an external file, use:
% (Replace '---' with actual filename)
% '\input{\modulefiles/tables/---.tex}' for tables
% '\includegraphics{\modulefiles/figures/---}' for figures
% '\subfile{\modulefiles/---}' for register files


% >>> M1: This module is in global project? <<<
% 'YES' - uncomment; 'NO' - comment out
\globaltrue

\ifglobal
    \documentclass[main\_\_CN.tex]{subfiles}

\else
    % Font "FandolSong-Regular" used by ctex does not have GSUB table.
    % It causes warnings that do not affect compilation and can be ignored.
    \PassOptionsToPackage{quiet}{fontspec}

    \documentclass[openany, 10pt]{book}
    % >>> M2: In ./chip-spec-settings/preamble-trm-module, also find
    %         settings for visibility of todo notes, watermark <<<

    \usepackage{preamble-trm-module}


    % Enable Chinese typesetting
    \usepackage{ctex}

    % >>> M3: Temporary labels in Chinese
    %         you can add your own labels there <<<
    \myexternaldocument{temp-labels-cn}

\fi %\ifglobal


\setcounter{secnumdepth}{4}


% >>> M5: If this module needs any updates in
% './chip-spec-settings/preamble-trm-module.sty'
% (include additional package, etc.), contact your module PM
% or create the file './chip-spec-settings/changelog.tex' make the
% necessary updates yourself and document them in this file <<<

\raggedright


\begin{document}


% Sets language into which pre-defined bits of text are translated
\selectlanguage{chinese}

% Do not delete this block
\ifglobal
% Add the following front matter if
% this module is outside of global project
\else
    \InputIfFileExists{readme}{}{}
    {\let\clearpage\relax \listoftodos}
    {\let\clearpage\relax \listofdonetodos}
    \tableofcontents
    %\listoftables
    %\listoffigures
\fi


%%%%%%%%%%%%%%%%%%%%%%%%%
%%% TRM Module Begins %%%
%%%%%%%%%%%%%%%%%%%%%%%%%

\hypertarget{vdma}{}
\chapter{VDMA 控制器 (VDMA)}
\label{mod:vdma}

%%%%%%%%%%%%%%%%%%%%
%%% Introduction %%%
%%%%%%%%%%%%%%%%%%%%

\section{概述}

DMA（直接内存访问）是指不依赖 CPU，在存储器和外设之间完成数据搬运。\chipname{} 上的 VDMA 控制器是一种通用 DMA，支持存储器到存储器、存储器到外设、外设到存储器的高速数据传输。VDMA 控制器以 AXI 作为总线接口，采用 AXI3 标准协议，带有两个 AXI 主机接口，支持用户动态地选择主机接口进行数据传输。

%%%%%%%%%%%%%%%%%%%%
%%% Terminology %%%
%%%%%%%%%%%%%%%%%%%%

\section{术语}

下列术语是在本章节的语境下定义的。

\begin{itemize}
    \item \textbf{AXI}：全称为 Advanced eXtensible Interface（高级可扩展接口），是 AMBA 3（Advanced Microcontroller Bus Architecture，高级微控制器总线架构）协议的一部分，为存储器和外设之间、存储器和存储器之间提供了高宽带、低延迟的数据传输。

    \item \textbf{外设}：指的是与 VDMA 之间有握手接口的系统模块，即表 \ref{tab:VDMA-hw-hs} 列出的外设。

    \item \textbf{存储器}：存储器总是为 DMA 传输做好准备，与 VDMA 之间没有握手接口。VDMA 支持的存储器包括 HP L2MEM、外部 flash 和外部 RAM。

    \item \textbf{源}：指的是源外设或源存储器，VDMA 通过 AXI 接口从源中读取数据，然后将数据存储在通道 FIFO 中。

    \item \textbf{源传输}：指的是源到 VDMA 的传输。

    \item \textbf{目标}：指的是目标外设或目标存储器，VDMA 将存储在通道 FIFO 中的数据写入目标。

    \item \textbf{目标传输}：指的是 VDMA 到目标的传输。

    \item \textbf{通道}：指的是源和目标之间的数据通路。

    \item \textbf{主机接口}：指的是 VDMA 的两个 AXI 主机接口。

    \item \textbf{从机接口}：指的是用于配置寄存器的 APB3 接口。

    \item \textbf{握手接口}：握手接口用于在 VDMA 和源或目标外设之间执行握手。握手接口用于请求、确认和控制 VDMA 传输。
    \item \textbf{流控制器}：决定 DMA 块传输的长度并终止传输的设备，可以是 VDMA、源或目标外设。存储器不能用作流控制器。

\end{itemize}

%%%%%%%%%%%%%%%%%%%%
%%% Feature List %%%
%%%%%%%%%%%%%%%%%%%%

\section{特性}

VDMA 具有以下特性：

\begin{itemize}
    \item 4 个通道，每个通道都是单向的，支持源到目标的数据传输
    \item 2 个 AXI 主机接口
    \item 支持与 MIPI DSI（显示串行接口）和 ISP（图像信号处理器）的握手
    \item 支持存储器与存储器之间、ISP 与存储器之间、MIPI DSI 与存储器之间的 DMA 传输
    \item DMA 传输层次有多个级别
    \item 每个通道的传输类型、传输长度、传输大小可配置
    \item 支持单块传输
    \item 支持基于连续地址、自动重新加载、影子寄存器和链表的多块传输
    \item 支持源传输和目标传输独立配置多块传输类型
    \item 通道禁用而不丢失数据
    \item 通道暂停、恢复和中止
    \item 可配置的通道优先级仲裁
    \item 支持使用 VDMA 或外设作为流控制器
    \item 外设的握手接口和通道之间的映射可配置
\end{itemize}


%%%%%%%%%%%%%%%%%%%%%%%%%%%%%%
%%% Architectural Overview %%%
%%%%%%%%%%%%%%%%%%%%%%%%%%%%%%

\section{结构概览}
\label{sec:vdma-arch-overview}

图 \ref{fig:VDMA-arch} 所示为 VDMA 的架构，以及 VDMA 与其他系统模块的连接。

\begin{figure}[H]
     \centering
     \includegraphics[width=0.9\linewidth]{\modulefiles/figures/VDMA-diagram.pdf}
     \caption{VDMA 架构}
     \label{fig:VDMA-arch}
 \end{figure}

从图中可以看出：

\begin{itemize}
    \item CPU 通过 APB3 从机接口 (APB3 Slave Interface) 访问 VDMA 的寄存器。
    \item VDMA 有 4 个通道（Channel 1 到 Channel 4）。每个通道对应一个 FIFO，深度为 64，宽度为 64 位。
    \item VDMA 有 2 个 AXI 主机接口（Master 1 和 Master 2）。
        \begin{itemize}
            \item Master 1 可以访问 ISP、MIPI DSI 和存储器，如图中蓝色箭头所示。
            \item Master 2 只能访问存储器，如图中红色箭头所示。
            \item Master 1 和 Master 2 各有一个仲裁器，对通道上的读请求和写请求进行仲裁。
        \end{itemize}
    \item VDMA 有 3 个握手接口。表 \ref{tab:VDMA-hw-hs} 为握手接口和外设之间的映射。
\end{itemize}

\begin{longtable}{ | p{6.8cm} | p{5.2cm} | }
\caption{握手接口与外设映射表}
\label{tab:VDMA-hw-hs} \\
\hline
\rowcolor{lightgray}
\textbf{握手接口编号} & \textbf{外设} \\ \hline
0 & MIPI DSI  \\ \hline
1 & ISP（ISP 到存储器） \\ \hline
2 & ISP（存储器到 ISP） \\ \hline
\end{longtable}

更多关于 ISP 和 MIPI DSI 的信息请参考章节 \ref{mod:isp} \textit{\nameref{mod:isp}}、\ref{mod:mipidsi} \textit{\nameref{mod:mipidsi}}。

%%%%%%%%%%%%%%%%%%%%%%%%%%%%%%
%%% Functional Description %%%
%%%%%%%%%%%%%%%%%%%%%%%%%%%%%%
\section{功能描述}

\subsection{传输层次}

传输最多分为 4 个级别：DMA 传输级别、块传输级别、transaction（事务传输）级别和 AXI 传输级别。

图 \ref{fig:VDMA-hier-non-mem} 所示为 VDMA 与外设之间传输的 4 个传输级别的层次结构。

图 \ref{fig:VDMA-hier-mem} 所示为 VDMA 与存储器之间传输的 3 个传输级别结构：DMA 传输、块传输和 AXI 传输。VDMA 与存储器之间的传输没有 transaction 级别，因为存储器始终假定为准备好进行数据传输。

以下关于 DMA 传输、块传输和 AXI 传输的描述均适用于存储器和外设。\textbf{transaction 级别的描述只适用于 VDMA 与外设之间的传输。}

\begin{itemize}
    \item \textbf{DMA 传输}：软件决定 DMA 传输中的块数。DMA 传输完成后，VDMA 会禁用通道并产生中断。然后可以对通道重新配置用以进行新的 DMA 传输。DMA 传输可能是单块传输，也可能是多块传输。

    \item \textbf{块传输}：块传输传输的是一个 VDMA 数据块，其数量由块长度表示，由流控制器决定。块传输有两种类型：单块传输和多块传输。由图 \ref{fig:VDMA-hier-non-mem} 可知，一个块传输可能由多个 transaction（事务传输）组成。

    \item \textbf{transaction（事务传输）}：一次外设请求是单次 transaction 或突发 transaction，表示需要一个 AXI 单次传输或多个 AXI 突发传输。

        \begin{itemize}
            \item 单次 transaction：长度始终为 1，单位为 \hyperref[fielddesc:DMACCHNSRCTRWIDTH]{DMAC\_CH\regindex{n}\_SRC\_TR\_WIDTH}。
            \item 突发 transaction：长度通过源的 \hyperref[fielddesc:DMACCHNSRCMSIZE]{DMAC\_CH\regindex{n}\_SRC\_MSIZE} 或目标的 \hyperref[fielddesc:DMACCHNDSTMSIZE]{DMAC\_CH\regindex{n}\_DST\_MSIZE} 被写到 VDMA 中，通常与源或目标外设中的 FIFO 大小有某种关系。具体请参考各外设章节。
        \end{itemize}
    \item \textbf{AXI 传输}：指的是 AXI 协议的传输，分为 AXI 单次传输和 AXI 突发传输。
    \end{itemize}

\begin{figure}[H]
     \centering   \includegraphics[width=\linewidth]{\modulefiles/figures/VDMA-Hier-Non-Mem-cn.pdf}
     \caption{外设的传输层次结构}
     \label{fig:VDMA-hier-non-mem}
 \end{figure}

\begin{figure}[H]
     \centering
     \includegraphics[width=0.9\linewidth]{\modulefiles/figures/VDMA-Hier-Mem-cn.pdf}
     \caption{存储器的传输层次结构}
     \label{fig:VDMA-hier-mem}
 \end{figure}

\subsection{仲裁机制}

仲裁机制对所有通道的读请求和写请求进行仲裁。

VDMA 采用优先级和公平平等 (fair-among-equals) 的仲裁机制。每个通道的优先级可在通道配置寄存器 \hyperref[regdesc:DMACCHNCFG1REG]{DMAC\_CH\regindex{n}\_CFG1\_REG} 中定义。

优先级的范围是 0（最低）到 3（最高），多个通道可以被赋予相同的优先级。

优先级和公平平等仲裁的仲裁机制工作原理如下：

\begin{itemize}
    \item 根据请求的优先级决定向哪个请求客户端发出授权信号。具有最高优先级的请求将被授予访问通道的权限。

    \item 两个或多个请求具有相同的优先级时，具有相同优先级的请求客户端将随机地获得权限。
\end{itemize}

\subsubsection{读仲裁器}

读仲裁器仲裁的是读请求，包括：

\begin{itemize}
    \item VDMA 获取源的数据或状态的读请求 (Source Data/Status Read Request)
    \item VDMA 获取目标的状态的读请求 (Destination Status Read Request)
    \item VDMA 获取链表描述符的读请求 (LLI Descriptor Read Request)
\end{itemize}

读仲裁器的示意图如图 \ref{fig:VDMA-read-arb} 所示。

\begin{figure}[H]
     \centering
     \includegraphics[width=1.0\linewidth]{\modulefiles/figures/VDMA-Read-Arb.pdf}
     \caption{VDMA 读仲裁器}
     \label{fig:VDMA-read-arb}
 \end{figure}

无论通道是否配置为最大优先级，获取链表描述符的读请求都将优先获得仲裁权限。

\subsubsection{写仲裁器}

写仲裁器仲裁的是写请求，包括：

\begin{itemize}
    \item 目标数据写请求 (Destination Data Write Request)
    \item 链表描述符回写请求 (LLI Descriptor Write-Back Request)
\end{itemize}

写仲裁器的示意图如图 \ref{fig:VDMA-write-arb} 所示。

\begin{figure}[H]
     \centering
     \includegraphics[width=1.0\linewidth]{\modulefiles/figures/VDMA-Write-Arb.pdf}
     \caption{VDMA 写仲裁器}
     \label{fig:VDMA-write-arb}
\end{figure}

在相同的通道中，目标数据写请求优先级高于链表描述符回写的优先级；如果不同通道具有相同的优先级，则目标数据写请求具有最高优先级。

\subsection{握手接口}

外设使用握手接口向 VDMA 表明它已准备好通过 AXI 总线来传输或接收数据。握手接口的操作取决于流控制器是外设还是 VDMA。

VDMA 在每次握手时执行一次单次 transaction 或突发 transaction。在某些情况下，仅使用突发 transaction 无法完成块传输。这种情况通常发生在块大小不是突发 transaction 长度的整数倍时。在这种情况下，块传输使用突发 transaction，直到块的剩余大小小于突发 transaction 中的数据量为止。此时，使用单次 transaction 完成块传输。

源的突发 transaction 大小由 \hyperref[fielddesc:DMACCHNSRCMSIZE]{DMAC\_CH\regindex{n}\_SRC\_MSIZE} 决定，目标的突发 transaction 大小由 \hyperref[fielddesc:DMACCHNDSTMSIZE]{DMAC\_CH\regindex{n}\_DST\_MSIZE} 决定，这也是用于每次握手的突发 transaction 的大小。\hyperref[fielddesc:DMACCHNSRCMSIZE]{DMAC\_CH\regindex{n}\_SRC\_MSIZE} 或 \hyperref[fielddesc:DMACCHNDSTMSIZE]{DMAC\_CH\regindex{n}\_DST\_MSIZE} 的值应与外设的配置一致。更多信息，请参考相关外设章节。

VDMA 与 MIPI DSI、ISP （ISP 到存储器）和 ISP（存储器到 ISP）之间有 3 个握手接口。表 \ref{tab:VDMA-hw-hs} 说明了握手接口和外设之间的映射关系。

\subsection{传输控制}

传输控制方便将数据从源搬运到目标。来自源的数据在被发送到目标之前暂时存储在通道的 FIFO 中。如果源和目标使用不同的传输大小进行数据传输，VDMA 将数据进行打包和解包，以适应 FIFO 的配置。

对于特定的 DMA 传输，传输类型和流控制配置由 \hyperref[fielddesc:DMACCHNTTFC]{DMAC\_CH\regindex{n}\_TT\_FC} 决定。

\subsubsection{单块传输}

如果一个 DMA 传输仅由单个块组成，软件可以将 \hyperref[fielddesc:DMACCHNSRCMULTBLKTYPE]{DMAC\_CH\regindex{n}\_SRC\_MULTBLK\_TYPE} 和 \hyperref[fielddesc:DMACCHNDSTMULTBLKTYPE]{DMAC\_CH\regindex{n}\_DST\_MULTBLK\_TYPE} 设置为 0，将 VDMA 配置为连续地址的单块传输。此时，VDMA 在完成 \hyperref[fielddesc:DMACCHNBLOCKTS]{DMAC\_CH\regindex{n}\_BLOCK\_TS} 定义的块长度传输时，将禁用通道。

\begin{tiplisting}
    单块传输其实只是多块传输的一种特例，两者并无严格区分。在本文档中，单块传输指的是只有一个块，多块传输至少包含两个块。
\end{tiplisting}

\subsubsection{多块传输}

如果一个 DMA 传输包含多个块，软件可以配置 \hyperref[fielddesc:DMACCHNSRCMULTBLKTYPE]{DMAC\_CH\regindex{n}\_SRC\_MULTBLK\_TYPE} 和 \hyperref[fielddesc:DMACCHNDSTMULTBLKTYPE]{DMAC\_CH\regindex{n}\_DST\_MULTBLK\_TYPE} 来选择多块传输类型。根据控制传输的寄存器（\hyperref[regdesc:DMACCHNSAR0REG]{DMAC\_CH\regindex{n}\_SAR0\_REG}、 \hyperref[regdesc:DMACCHNDAR0REG]{DMAC\_CH\regindex{n}\_DAR0\_REG}、 \hyperref[fielddesc:DMACCHNBLOCKTS]{DMAC\_CH\regindex{n}\_BLOCK\_TS}、 \hyperref[regdesc:DMACCHNCTL0REG]{DMAC\_CH\regindex{n}\_CTL0/1\_REG}）的更新方式不同，共有以下四种多块传输类型：
\begin{enumerate}
    \item 连续地址
    \item 自动重新加载，寄存器的值从初始值重新加载
    \item 影子寄存器，寄存器的值从影子寄存器加载
    \item 链表，寄存器的值从下一个链表加载
\end{enumerate}
表 \ref{tab:VDMA-reg-upd-methd} 列出了寄存器更新方法的所有可能组合。

\begin{small}
\begin{longtable}{ | L{2.3cm} |L{2.3cm}| L{1.8cm} | L{1.8cm} | L{2.2cm} | L{1.8cm} | L{2cm} | }

\caption{多块传输寄存器更新方法}

\label{tab:VDMA-reg-upd-methd} \\
\hline
\rowcolor{lightgray}

\multicolumn{2}{|c|}{\textbf{\hyperref[regdesc:DMACCHNCFG0REG]{DMAC\_CH\regindex{n}\_CFG0\_REG}}} & \multicolumn{4}{|c|}{\textbf{寄存器更新方法}} & \\ \cline{1-6}
\rowcolor{lightgray}

\textbf{\hyperref[fielddesc:DMACCHNSRCMULTBLKTYPE]{\makecell[l]{DMAC\_CH\regindex{n}\_\\SRC\_MULTBLK\_\\TYPE}}} & \textbf{\hyperref[fielddesc:DMACCHNDSTMULTBLKTYPE]{\makecell[l]{DMAC\_CH\regindex{n}\_\\DST\_MULTBLK\_\\TYPE}}} & \textbf{\hyperref[regdesc:DMACCHNSAR0REG]{\makecell[l]{DMAC\_CH\regindex{n}\_\\SAR0\_REG}}} & \textbf{\hyperref[regdesc:DMACCHNDAR0REG]{\makecell[l]{DMAC\_CH\regindex{n}\_\\DAR0\_REG}}} & \textbf{\hyperref[regdesc:DMACCHNCTL0REG]{\makecell[l]{DMAC\_CH\regindex{n}\_\\CTL0\_REG}}/\hyperref[regdesc:DMACCHNCTL1REG]{CTL1\_REG}} & \textbf{\hyperref[fielddesc:DMACCHNBLOCKTS]{\makecell[l]{DMAC\_CH\regindex{n}\_\\BLOCK\_TS}}} & \textbf{说明} \\ \hline
00 & 00 & 连续地址 & 连续地址 & 不更新 & 不更新 & 单块或多块的最后一块  \\ \hline
00 & 01 & 连续地址 & 从初始值重新加载 & 从初始值重新加载 & 从初始值重新加载 & -  \\ \hline
00 & 10 & 连续地址 & 从影子寄存器加载 & 从影子寄存器加载 & 从影子寄存器加载 & -  \\ \hline
01 & 00 & 从初始值重新加载 & 连续地址 & 从初始值重新加载 & 从初始值重新加载 & -  \\ \hline
01 & 01 & 从初始值重新加载 & 从初始值重新加载 & 从初始值重新加载 & 从初始值重新加载 & -  \\ \hline
01 & 10 & 从初始值重新加载 & 从影子寄存器加载 & 从影子寄存器加载 & 从影子寄存器加载 & -  \\ \hline
10 & 00 & 从影子寄存器加载 & 连续地址 & 从影子寄存器加载 & 从影子寄存器加载 & -  \\ \hline
10 & 01 & 从影子寄存器加载 & 从初始值重新加载 & 从影子寄存器加载 & 从影子寄存器加载 & -  \\ \hline
10 & 10 & 从影子寄存器加载 & 从影子寄存器加载 & 从影子寄存器加载 & 从影子寄存器加载 & -  \\ \hline
00 & 11 & 连续地址 & 从下一个链表加载 & 从下一个链表加载 & 从下一个链表加载 & -  \\ \hline
01 & 11 & 从初始值重新加载 & 从下一个链表加载 & 从下一个链表加载 & 从下一个链表加载 & -  \\ \hline
11 & 00 & 从下一个链表加载 & 连续地址 & 从下一个链表加载 & 从下一个链表加载 & -  \\ \hline
11 & 01 & 从下一个链表加载 & 从初始值重新加载 & 从下一个链表加载 & 从下一个链表加载 & -  \\ \hline
11 & 11 & 从下一个链表加载 & 从下一个链表加载 & 从下一个链表加载 & 从下一个链表加载 & -  \\ \hline
10 & 11 & - & - & - & - & 无效。生成中断* \\ \hline
11 & 10 & - & - & - & - & 无效。生成中断*\\ \hline

\end{longtable}
\end{small}
\vspace{-1em}
*说明：生成的中断为 \hyperref[int:DMACCHNSLVIFMULTIBLKTYPEERRINT]{DMAC\_CH\regindex{n}\_SLVIF\_MULTIBLKTYPE\_ERR\_INT}。

\subparagraph{连续地址}

在这种情况下，每个连续块的地址从前一个块的结束地址开始。配置 \hyperref[fielddesc:DMACCHNSRCMULTBLKTYPE]{DMAC\_CH\regindex{n}\_SRC\_MULTBLK\_TYPE} 或 \hyperref[fielddesc:DMACCHNDSTMULTBLKTYPE]{DMAC\_CH\regindex{n}\_DST\_MULTBLK\_TYPE} 为 0 使源地址或目标地址在块之间连续。

在多块传输的情况下，\hyperref[fielddesc:DMACCHNSRCMULTBLKTYPE]{DMAC\_CH\regindex{n}\_SRC\_MULTBLK\_TYPE} 和 \hyperref[fielddesc:DMACCHNDSTMULTBLKTYPE]{DMAC\_CH\regindex{n}\_DST\_MULTBLK\_TYPE} 的值不能同时被选择为连续地址，即不能同时为 0。若某一块传输过程中 \hyperref[fielddesc:DMACCHNSRCMULTBLKTYPE]{DMAC\_CH\regindex{n}\_SRC\_MULTBLK\_TYPE} 和 \hyperref[fielddesc:DMACCHNDSTMULTBLKTYPE]{DMAC\_CH\regindex{n}\_DST\_MULTBLK\_TYPE} 都被配置成了 0，那么当前传输就会被认为是最后一个块。

在多块传输中，如果需要 \hyperref[regdesc:DMACCHNSAR0REG]{DMAC\_CH\regindex{n}\_SAR0\_REG} 和 \hyperref[regdesc:DMACCHNDAR0REG]{DMAC\_CH\regindex{n}\_DAR0\_REG} 在块之间地址连续，可以通过使用链表或影子寄存器间接实现此功能。

\begin{tiplisting}
如果启用了基于连续地址的多块传输（即 \hyperref[fielddesc:DMACCHNSRCMULTBLKTYPE]{DMAC\_CH\regindex{n}\_SRC\_MULTBLK\_TYPE} 或 \hyperref[fielddesc:DMACCHNDSTMULTBLKTYPE]{DMAC\_CH\regindex{n}\_DST\_MULTBLK\_TYPE} 为 0），则在完整的 DMA 传输中必须至少有两个块，否则会导致不可预测的行为。

\end{tiplisting}

\subparagraph{从初始值重新加载}

在这种情况下，\hyperref[regdesc:DMACCHNSAR0REG]{DMAC\_CH\regindex{n}\_SAR0\_REG} 和 \hyperref[regdesc:DMACCHNDAR0REG]{DMAC\_CH\regindex{n}\_DAR0\_REG} 通道寄存器在新的块传输开始时从它们的初始值重新加载。
VDMA 不会进行下一个块传输，直到软件清除相应通道的块传输完成中断，即将 \hyperref[fielddesc:DMACCHNCLEARSRCTRANSCOMPINTSTAT]{DMAC\_CH\regindex{n}\_CLEAR\_SRC\_TRANSCOMP\_INTSTAT} 和 \hyperref[fielddesc:DMACCHNCLEARDSTTRANSCOMPINTSTAT]{DMAC\_CH\regindex{n}\_CLEAR\_DST\_TRANSCOMP\_INTSTAT} 置为 1。

\begin{tiplisting}
    如果启用了基于自动重新加载的多块传输，则在完整的 DMA 传输中必须至少有两个块，否则会导致不可预测的行为。
\end{tiplisting}

\subparagraph{从影子寄存器加载}

在这种情况下，\hyperref[regdesc:DMACCHNSAR0REG]{DMAC\_CH\regindex{n}\_SAR0\_REG}、\hyperref[regdesc:DMACCHNDAR0REG]{DMAC\_CH\regindex{n}\_DAR0\_REG}、\hyperref[fielddesc:DMACCHNBLOCKTS]{DMAC\_CH\regindex{n}\_BLOCK\_TS}、\hyperref[regdesc:DMACCHNCTL0REG]{DMAC\_CH\regindex{n}\_CTL0\_REG} 和 \hyperref[regdesc:DMACCHNCTL1REG]{DMAC\_CH\regindex{n}\_CTL1\_REG} 在每个块完成时与它们对应的影子寄存器进行同步。影子寄存器的值将用于后续的块。

软件将对应块的配置写到相应的影子寄存器。在完成上一个块传输时，VDMA 将影子寄存器的内容复制到 \hyperref[regdesc:DMACCHNSAR0REG]{DMAC\_CH\regindex{n}\_SAR0\_REG}、\hyperref[regdesc:DMACCHNDAR0REG]{DMAC\_CH\regindex{n}\_DAR0\_REG}、\hyperref[fielddesc:DMACCHNBLOCKTS]{DMAC\_CH\regindex{n}\_BLOCK\_TS}、\hyperref[regdesc:DMACCHNCTL0REG]{DMAC\_CH\regindex{n}\_CTL0\_REG} 和 \hyperref[regdesc:DMACCHNCTL1REG]{DMAC\_CH\regindex{n}\_CTL1\_REG} 然后再开始新的块传输。

对 \hyperref[regdesc:DMACCHNSAR0REG]{DMAC\_CH\regindex{n}\_SAR0\_REG}、\hyperref[regdesc:DMACCHNDAR0REG]{DMAC\_CH\regindex{n}\_DAR0\_REG}、\hyperref[fielddesc:DMACCHNBLOCKTS]{DMAC\_CH\regindex{n}\_BLOCK\_TS}、\hyperref[regdesc:DMACCHNCTL0REG]{DMAC\_CH\regindex{n}\_CTL0\_REG} 和 \hyperref[regdesc:DMACCHNCTL1REG]{DMAC\_CH\regindex{n}\_CTL1\_REG} 的读操作总是返回与当前块传输对应的数据，而不是对应下一个块的影子寄存器的内容。

\hyperref[fielddesc:DMACCHNSHADOWREGORLLIVALID]{DMAC\_CH\regindex{n}\_SHADOWREG\_OR\_LLI\_VALID} 指示影子寄存器的内容是否有效（1 为有效，0 为无效）。如果在影子寄存器获取阶段，此位的值为 0，VDMA 将忽略影子寄存器的内容并生成 \hyperref[int:DMACCHNSHADOWREGORLLIINVALIDERRINT]{DMAC\_CH\regindex{n}\_SHADOWREG\_OR\_LLI\_INVALID\_ERR\_INT} 中断。VDMA 等待软件将 1 写入 \hyperref[fielddesc:DMACCHNBLKTFRRESUMEREQ]{DMAC\_CH\regindex{n}\_BLK\_TFR\_RESUMEREQ} 以指示影子寄存器可用，然后再尝试读取影子寄存器以继续下一个块传输。
\begin{tiplisting}
    在从影子寄存器加载的多块传输下，\hyperref[regdesc:DMACCHNSAR0REG]{DMAC\_CH\regindex{n}\_SAR0\_REG}、\hyperref[regdesc:DMACCHNDAR0REG]{DMAC\_CH\regindex{n}\_DAR0\_REG}、\hyperref[fielddesc:DMACCHNBLOCKTS]{DMAC\_CH\regindex{n}\_BLOCK\_TS}、\\\hyperref[regdesc:DMACCHNCTL0REG]{DMAC\_CH\regindex{n}\_CTL0\_REG}、\hyperref[regdesc:DMACCHNCTL1REG]{DMAC\_CH\regindex{n}\_CTL1\_REG} 寄存器里的实际值是当前传输对应的配置。软件写这些寄存器其实是写到对应的影子寄存器里面。
\end{tiplisting}

\subparagraph{链表}

在这种情况下，VDMA 会从存储器中获取某个块的块描述符，用来更新 \hyperref[regdesc:DMACCHNSAR0REG]{DMAC\_CH\regindex{n}\_SAR0\_REG}、\hyperref[regdesc:DMACCHNDAR0REG]{DMAC\_CH\regindex{n}\_DAR0\_REG}、\hyperref[fielddesc:DMACCHNBLOCKTS]{DMAC\_CH\regindex{n}\_BLOCK\_TS}、\hyperref[regdesc:DMACCHNCTL0REG]{DMAC\_CH\regindex{n}\_CTL0\_REG} 和 \hyperref[regdesc:DMACCHNCTL1REG]{DMAC\_CH\regindex{n}\_CTL1\_REG}，然后再发起块传输。这个过程被称为 LLI（链表项）更新。VDMA 的块链 (block chaining) 特性使用链表指针寄存器 \hyperref[fielddesc:DMACCHNLOC0]{DMAC\_CH\regindex{n}\_LOC0} 来存储下一个链表项在存储器中的地址，以此实现数据传输的无缝连续。
图 \ref{fig:VDMA-lli} 显示了 VDMA 链表项。

\begin{figure}[H]
    \centering
    \begin{bytefield}[bitwidth=1.2em, boxformatting={\centering\small}, leftcurly=., leftcurlyspace=0pt]{32}
        \bitheader[bitformatting={\scriptsize}]{0,31} \\
        \bitbox{32}{Reserved}\\
        \bitbox{32}{Reserved}\\
        \bitbox{32}{CHn\_LLP\_STATUS[63:32]}\\
        \bitbox{32}{CHn\_LLP\_STATUS[31:0]}\\
        \bitbox{32}{Write Back for CHn\_DSTAT}\\
        \bitbox{32}{Write Back for CHn\_SSTAT}\\
        \bitbox{32}{CHn\_CTL1[31:0]}\\
        \bitbox{32}{CHn\_CTL0[31:0]}\\
        \bitbox{32}{Reserved}\\
            \bitbox{26}{CHn\_LLP0[31:6]} & \bitbox{6}{Reserved[5:0]}\\
        \bitbox{32}{Reserved}\\
        \bitbox{32}{CHn\_BLOCK\_TS[31:0]}\\
        \bitbox{32}{Reserved}\\
        \bitbox{32}{CHn\_DAR[31:0]}\\
        \bitbox{32}{Reserved}\\
        \bitbox{32}{CHn\_SAR[31:0]}
    \end{bytefield}
    \caption{VDMA 链表项}
    \label{fig:VDMA-lli}
\end{figure}

VDMA 具有动态扩展链表的特性，不需要提前在系统内存中创建整个链表。要实现此功能，可以将 \hyperref[fielddesc:DMACCHNSHADOWREGORLLIVALID]{DMAC\_CH\regindex{n}\_SHADOWREG\_OR\_LLI\_VALID} 配置为 1，如果链表项是最后一个则 \hyperref[fielddesc:DMACCHNSHADOWREGORLLILAST]{DMAC\_CH\regindex{n}\_SHADOWREG\_OR\_LLI\_LAST} 也应配置为 1。

对于基于链表的多块传输，\hyperref[fielddesc:DMACCHNSHADOWREGORLLIVALID]{DMAC\_CH\regindex{n}\_SHADOWREG\_OR\_LLI\_VALID} 用来指示从存储器中获取的链表项是否有效（1 为有效，0 为无效）。如果 LLI 无效，VDMA 将丢弃 LLI 并生成 \hyperref[int:DMACCHNSHADOWREGORLLIINVALIDERRINT]{DMAC\_CH\regindex{n}\_SHADOWREG\_OR\_LLI\_INVALID\_ERR\_INT} 中断。此中断会导致 VDMA 停止相应的通道，等待软件将 \hyperref[fielddesc:DMACCHNBLKTFRRESUMEREQ]{DMAC\_CH\regindex{n}\_BLK\_TFR\_RESUMEREQ} 写为 1 以指示 LLI 可用，然后再尝试读取 LLI。

\begin{tiplisting}
    对于预取的 LLI，如果 \hyperref[fielddesc:DMACCHNSHADOWREGORLLIVALID]{DMAC\_CH\regindex{n}\_SHADOWREG\_OR\_LLI\_VALID} 为 0，则不会生成\\ \hyperref[fielddesc:DMACCHNSHADOWREGORLLIINVALIDERRINTSTAT]{DMAC\_CH\regindex{n}\_SHADOWREG\_OR\_LLI\_INVALID\_ERR\_INTSTAT} 中断。VDMA 在完成当前块传输后，再次尝试获取 LLI，只有当 \hyperref[fielddesc:DMACCHNSHADOWREGORLLIVALID]{DMAC\_CH\regindex{n}\_SHADOWREG\_OR\_LLI\_VALID} 仍然为 0 时，才会生成\\ \hyperref[fielddesc:DMACCHNSHADOWREGORLLIINVALIDERRINTSTAT]{DMAC\_CH\regindex{n}\_SHADOWREG\_OR\_LLI\_INVALID\_ERR\_INTSTAT} 中断。
\end{tiplisting}

下列 LLI 字段由软件准备：

\begin{itemize}
    \item CH\regindex{n}\_SAR：字段描述同 \hyperref[regdesc:DMACCHNSAR0REG]{DMAC\_CH\regindex{n}\_SAR0\_REG}
    \item CH\regindex{n}\_DAR：字段描述同 \hyperref[regdesc:DMACCHNDAR0REG]{DMAC\_CH\regindex{n}\_DAR0\_REG}
    \item CH\regindex{n}\_BLOCK\_TS：字段描述同 \hyperref[fielddesc:DMACCHNBLOCKTS]{DMAC\_CH\regindex{n}\_BLOCK\_TS}
    \item CH\regindex{n}\_LLP0[31:6]：字段描述同 \hyperref[regdesc:DMACCHNLLP0REG]{DMAC\_CH\regindex{n}\_LLP0\_REG}[31:6]
    \item CH\regindex{n}\_CTL0：字段描述同 \hyperref[regdesc:DMACCHNCTL0REG]{DMAC\_CH\regindex{n}\_CTL0\_REG}
    \item CH\regindex{n}\_CTL1：字段描述同 \hyperref[regdesc:DMACCHNCTL1REG]{DMAC\_CH\regindex{n}\_CTL1\_REG}
\end{itemize}

如果启用了状态回写选项，VDMA 将 CH\regindex{n}\_CTL0、CH\regindex{n}\_CTL1、CH\regindex{n}\_LLP\_STATUS、CH\regindex{n}\_SSTAT 和 CH\regindex{n}\_DSTAT 信息回写到为该字段定义的位置，即从地址[CH\regindex{n}\_LLP] + 0x20 到 [CH\regindex{n}\_LLP] + 0x34。CH\regindex{n}\_SSTAT 和 CH\regindex{n}\_DSTAT 的回写可以通过配置 \hyperref[fielddesc:DMACCHNSRCSTATEN]{DMAC\_CH\regindex{n}\_SRC\_STAT\_EN} 和 \hyperref[fielddesc:DMACCHNDSTSTATEN]{DMAC\_CH\regindex{n}\_DST\_STAT\_EN} 单独启用或禁用。

\begin{tiplisting}
\vspace{-2em}
\begin{itemize}
    \item 如果 CH\regindex{n}\_SSTAT 和 CH\regindex{n}\_DSTAT 回写没有启用，不要使用 PSRAM 来存储 LLI。
    \item LLI 回写操作后，\hyperref[fielddesc:DMACCHNSHADOWREGORLLIVALID]{DMAC\_CH\regindex{n}\_SHADOWREG\_OR\_LLI\_VALID} 将为 0。
\end{itemize}

\end{tiplisting}

图 \ref{fig:VDMA-llp} 显示了 CH\regindex{n}\_LLP\_STATUS 回写字段：

\begin{figure}[H]

    \begin{bytefield}[bitwidth=0.75em, bitheight=3em, boxformatting={\centering\small}, leftcurly=., leftcurlyspace=0pt]{64}
        \bitheader[bitformatting={\scriptsize}]{0,63} \\
        \bitbox{2}{\textbf{*}} & \bitbox{15}{Reserved[61:47]} & \bitbox{15}{Data Left in Channel FIFO (CHn\_Status[46:32])} & \bitbox{10}{Reserved[31:22]} & \bitbox{22}{Completed Block Transfer Size (CHn\_Status[21:0])}
    \end{bytefield}
说明：\textbf{*} 代表 Status Indication (CHn\_IntStatus[1:0])
    \caption{CH\regindex{n}\_LLP\_STATUS 回写字段}
    \label{fig:VDMA-llp}
\end{figure}

CH\regindex{n}\_LLP\_STATUS[63] 和 CH\regindex{n}\_LLP\_STATUS[62] 分别表示 DMA\_TFR\_DONE 和 BLOCK\_TFR\_DONE 状态。这两个字段是 LLI 回写过程中最后更新的字段。软件应该确保 BLOCK\_TFR\_DONE 为 1 再使用 CH\regindex{n}\_SSTAT 和 CH\regindex{n}\_DSTAT 信息。 DMA\_TFR\_DONE 和 BLOCK\_TFR\_DONE 在传输最后一个块后会被置为 1。

CH\regindex{n}\_STATUS[21:0] 对应 \hyperref[regdesc:DMACCHNSTATUS0REG]{DMAC\_CH\regindex{n}\_CMPLTD\_BLK\_TFR\_SIZE}。CH\regindex{n}\_STATUS[46:32] 对应 \hyperref[regdesc:DMACCHNSTATUS1REG]{DMAC\_CH\regindex{n}\_DATA\_LEFT\_IN\_FIFO[14:0]}。
\subparagraph{块间传输暂停}

在每个块传输结束时，如果满足以下条件，则产生块传输完成中断 \hyperref[int:DMACCHNBLOCKTFRDONEINT]{DMAC\_CH\regindex{n}\_BLOCK\_TFR\_DONE\_INT}：
\begin{itemize}
    \item 使能全局中断（\hyperref[fielddesc:DMACINTEN]{DMAC\_INT\_EN} 置 1）
    \item 使能通道块传输完成中断（\hyperref[fielddesc:DMACCHNENABLEBLOCKTFRDONEINTSTAT] {DMAC\_CH\regindex{n}\_ENABLE\_BLOCK\_TFR\_DONE\_INTSTAT} 置 1，\hyperref[fielddesc:DMACCHNENABLEBLOCKTFRDONEINTSIGNAL]{DMAC\_CH\regindex{n}\_ENABLE\_BLOCK\_TFR\_DONE\_INTSIGNAL} 置 1，\hyperref[fielddesc:DMACCHNIOCBLKTFR]{DMAC\_CH\regindex{n}\_IOC\_BLKTFR} 置 1）
\end{itemize}

对于基于连续地址和自动重新加载的多块传输（源和目标外设都不使用基于影子寄存器或链表的多块传输），DMA 传输会在中断  \hyperref[int:DMACCHNBLOCKTFRDONEINT]{DMAC\_CH\regindex{n}\_BLOCK\_TFR\_DONE\_INT} 置 1 后自动停止。VDMA 不会继续进行下一个块传输，直到软件将 \hyperref[fielddesc:DMACCHNCLEARBLOCKTFRDONEINTSTAT]{DMAC\_CH\regindex{n}\_CLEAR\_BLOCK\_TFR\_DONE\_INTSTAT} 置 1 清除中断。

块之间的通道暂停用于确保倒数第二个块的块传输完成 ISR（中断服务例程）在最后一个块传输开始之前得到响应。这确保 ISR 在最后一个块传输完成之前清零 \hyperref[fielddesc:DMACCHNSRCMULTBLKTYPE]{DMAC\_CH\regindex{n}\_SRC\_MULTBLK\_TYPE} 和 \hyperref[fielddesc:DMACCHNDSTMULTBLKTYPE]{DMAC\_CH\regindex{n}\_DST\_MULTBLK\_TYPE}。

\subparagraph{多块传输结束}

如果源外设或目标外设使用基于影子寄存器或链表的多块传输，则指示位 \hyperref[fielddesc:DMACCHNSHADOWREGORLLILAST]{DMAC\_CH\regindex{n}\_SHADOWREG\_OR\_LLI\_LAST} 表明当前块是否是最后一个。如果 \hyperref[fielddesc:DMACCHNSHADOWREGORLLILAST]{DMAC\_CH\regindex{n}\_SHADOWREG\_OR\_LLI\_LAST} 为 1，VDMA 认为当前块是最后一块，并在当前块传输结束时完成 DMA 传输操作。

对于基于连续地址和自动重新加载的多块传输（源和目标外设都不使用基于影子寄存器或链表的多块传输），如果对应的多块类型寄存器 \hyperref[fielddesc:DMACCHNSRCMULTBLKTYPE]{DMAC\_CH\regindex{n}\_SRC\_MULTBLK\_TYPE} 和/或 \hyperref[fielddesc:DMACCHNDSTMULTBLKTYPE]{DMAC\_CH\regindex{n}\_DST\_MULTBLK\_TYPE} 在块传输结束时为 0，VDMA 认为前一个块为最后一个块，并完成 DMA 传输操作。

\subsection{流控制器}

DMA 块传输的长度由流控制器决定。流控制器可以是 VDMA、源或目标外设，存储器不能用作流控制器。

如果在启用通道之前知道块的长度，那么应该将 VDMA 配置为流控制器。如果在启用通道之前不知道块的长度，源或目标外设应该被配置为流控制器，并由其决定何时终止 DMA 块传输。

\subsection{通道挂起和恢复}

软件可以在 DMA 传输时挂起和恢复通道。挂起通道的操作如下：
\begin{enumerate}
    \item 软件将 \hyperref[fielddesc:DMACCH1SUSP]{DMAC\_CH\regindex{n}\_SUSP} 置为 1。
    \item VDMA 在完成源已发起的所有 AXI 传输后，暂停来自源的所有传输。
    \item VDMA 将 \hyperref[int:DMACCHNCHSRCSUSPENDEDINT]{DMAC\_CH\regindex{n}\_CH\_SRC\_SUSPENDED\_INT} 置为 1 表示源传输被挂起并产生一个挂起中断。

    \item VDMA 将通道 FIFO 里的数据全部传输给目标。当 \hyperref[fielddesc:DMACCHNSRCTRWIDTH]{DMAC\_CH\regindex{n}\_SRC\_TR\_WIDTH} < \hyperref[fielddesc:DMACCHNDSTTRWIDTH]{DMAC\_CH\regindex{n}\_DST\_TR\_WIDTH} 并且 \hyperref[fielddesc:DMACCH1SUSP]{DMAC\_CH\regindex{n}\_SUSP} 为 1 时，通道 FIFO 中可能还有数据，但不足以形成一次 \hyperref[fielddesc:DMACCHNDSTTRWIDTH]{DMAC\_CH\regindex{n}\_DST\_TR\_WIDTH} 宽度的单次传输。当通道恢复后，FIFO 中剩余的数据会被传输到目标。

    \item VDMA 将 \hyperref[int:DMACCHNCHSUSPENDEDINT]{DMAC\_CH\regindex{n}\_CH\_SUSPENDED\_INT} 置为 1 表示通道被挂起，并产生中断。

\end{enumerate}

通道被挂起后，软件可以将 \hyperref[fielddesc:DMACCH1SUSP]{DMAC\_CH\regindex{n}\_SUSP} 置 0，退出通道挂起模式。然后 VDMA 将恢复 DMA 传输。

\subsection{通道禁用}

在正常情况下，软件通过将 \hyperref[fielddesc:DMACCH1EN]{DMAC\_CH\regindex{n}\_EN} 置为 1 来启用一个通道，VDMA 在通道传输结束时清零 \hyperref[fielddesc:DMACCH1EN]{DMAC\_CH\regindex{n}\_EN}。软件可以在传输完成之前禁用通道。

\subsubsection{传输完成前禁用被挂起的通道}

要在传输完成前禁用被挂起的通道：

\begin{itemize}
    \item 在 VDMA 将 \hyperref[int:DMACCHNCHSUSPENDEDINT]{DMAC\_CH\regindex{n}\_CH\_SUSPENDED\_INT} 置 1 后将 \hyperref[fielddesc:DMACCH1EN]{DMAC\_CH\regindex{n}\_EN} 置为 0。当 \hyperref[fielddesc:DMACCHNSRCTRWIDTH]{DMAC\_CH\regindex{n}\_SRC\_TR\_WIDTH} < \hyperref[fielddesc:DMACCHNDSTTRWIDTH]{DMAC\_CH\regindex{n}\_DST\_TR\_WIDTH} 并且 \hyperref[fielddesc:DMACCH1SUSP]{DMAC\_CH\regindex{n}\_SUSP} 为 1 时，通道 FIFO 中可能还有数据，但不足以形成一次 \hyperref[fielddesc:DMACCHNDSTTRWIDTH]{DMAC\_CH\regindex{n}\_DST\_TR\_WIDTH} 宽度的单次传输。在这种情况下，FIFO 中剩余的数据不会被传输到目标，会丢失。

    \item VDMA 将 \hyperref[int:DMACCHNCHDISABLEDINT]{DMAC\_CH\regindex{n}\_CH\_DISABLED\_INT} 置为 1 表示通道被禁用，并产生中断。

\end{itemize}

\subsubsection{传输完成前禁用未被挂起的通道}

要在传输完成前禁用未被挂起的通道：

\begin{itemize}
    \item 软件在 DMA 传输时将 \hyperref[fielddesc:DMACCH1EN]{DMAC\_CH\regindex{n}\_EN} 置 0。
    \item VDMA 在完成源已发起的所有 AXI 传输后，停止来自源的所有传输。
    \item VDMA 将通道 FIFO 里的数据全部传输给目标。当 \hyperref[fielddesc:DMACCHNSRCTRWIDTH]{DMAC\_CH\regindex{n}\_SRC\_TR\_WIDTH} < \hyperref[fielddesc:DMACCHNDSTTRWIDTH]{DMAC\_CH\regindex{n}\_DST\_TR\_WIDTH} 并且 \hyperref[fielddesc:DMACCH1EN]{DMAC\_CH\regindex{n}\_EN} 为 0 时，通道 FIFO 中可能还有数据，但不足以形成一次 \hyperref[fielddesc:DMACCHNDSTTRWIDTH]{DMAC\_CH\regindex{n}\_DST\_TR\_WIDTH} 宽度的单次传输。在这种情况下，FIFO 中剩余的数据不会被传输到目标，会丢失。
    \item VDMA 将 \hyperref[int:DMACCHNCHDISABLEDINT]{DMAC\_CH\regindex{n}\_CH\_DISABLED\_INT} 置为 1 表示通道被禁用，并产生中断。

\end{itemize}

\begin{tiplisting}
    一旦软件将 \hyperref[fielddesc:DMACCH1EN]{DMAC\_CH\regindex{n}\_EN} 置 0，发起了通道禁用程序，用户无法在 \hyperref[int:DMACCHNCHDISABLEDINT]{DMAC\_CH\regindex{n}\_CH\_DISABLED\_INT} 置 1 之前将 \hyperref[fielddesc:DMACCH1EN]{DMAC\_CH\regindex{n}\_EN} 置 1 来重新使能通道。VDMA 会忽略这个写操作。
\end{tiplisting}

\subsection{低功耗技术}

VDMA 支持低功耗技术，这是一种智能低功耗逻辑，可以在各种场景下自动感知多个 DMA 通道、从机总线接口和 AXI 主机接口通道上的空闲时间，然后切断这些模块的时钟，从而降低总体功耗。低功耗逻辑反应迅速，可以在监测到任何活动时向这些模块提供时钟，而不会影响性能。

\hyperref[regdesc:DMACLOWPOWERCFG0REG]{DMAC\_LOWPOWER\_CFG0\_REG} 可以用来启用低功耗并控制低功耗计数器。

\subsubsection{DMA 通道低功耗技术}

DMA 通道占用了 VDMA 上的大部分面积，比其他模块消耗更多的功耗，因此更加需要优化其功耗。

DMA 通道低功耗技术使用 DMA 通道低功耗状态机监测通道上的活动。每当监测到 DMA 通道的低功耗条件时，DMA 通道低功耗延迟计数器就开始启动。状态机等待直到计数器到期，然后 DMA 通道通过关闭 VDMA 的时钟 AXI\_DMAC\_CORE\_CLK 进入低功耗状态。DMA 通道一直保持这个状态，直到在通道上检测到任何活动。在某些情况下，如果在延迟计数器运行时监测到活动，则计数器将重置并停止计数，直到再次检测到低功耗条件。
\subsubsection{从机总线接口上的低功耗技术}

从机总线接口单元 (SBIU) 通过 APB3 从机接口访问 VDMA 的内部寄存器。低功耗技术可以优化 SBIU 的功耗。

从机总线接口上的低功耗技术使用 SBIU 状态机监测从机总线接口上的活动。一旦检测到低功耗条件（即没有 APB 读写传输正在进行并且 SBIU 状态机处于空闲状态），SBIU 低功耗延迟计数器将开始计数。当延迟计数器到期时，SBIU 通过关闭 PCLK 或 AXI\_DMAC\_CORE\_CLK（根据时钟模式配置）进入低功耗状态。SBIU 一直保持这种状态，直到在从机总线接口上检测到读或写传输。在某些情况下，如果在延迟计数器运行时监测到传输，则计数器将重置并停止计数，直到再次检测到低功耗条件。

如果在 SBIU 处于低功耗状态时检测到读或写传输，则低功耗逻辑会立即响应，打开 PCLK 或 AXI\_DMAC\_CORE\_CLK 从而退出低功耗状态。

\subsubsection{AXI 主机接口通道上的低功耗技术}

AXI 主机接口通过 AXI 总线在存储器和外设之间传输数据。

每当检测到 AXI 主机接口通道上的低功耗条件时，AXI 主机接口通道低功耗延迟计数器开始计时。AXI 低功耗逻辑等待直到延迟计数器到期，然后 AXI 主机接口通道关闭 AXI 主机接口的时钟和 AXI\_DMAC\_CORE\_CLK 进入低功耗状态。AXI 主机接口通道一直保持这个状态，直到在通道上检测到任何活动。在某些情况下，如果在延迟计数器运行时监测到活动，则计数器将重置并停止计数，直到再次检测到低功耗条件。

\subsubsection{全局低功耗技术}

全局低功耗技术根据 DMA 通道、从机总线接口和 AXI 主机接口通道的空闲状态实现低功耗选项。如果上述模块都指示 VDMA 处于空闲状态，则全局低功耗延迟计数器将启动。

全局低功耗逻辑等待直到低功耗延迟计数器到期，然后 VDMA 关闭所有模块的时钟来进入低功耗状态。VDMA 一致保持这个状态，直到在任意模块上检测到任何活动。在某些情况下，如果在延迟计数器运行时监测到活动，则计数器将重置并停止计数，直到再次检测到低功耗条件。

%%%%%%%%%%%%%%%%%%%%%
%%% Interrupts %%%
%%%%%%%%%%%%%%%%%%%%%

\section{中断}

\chipname{} 的 VDMA 可以生成 AXI\_DMAC\_INTR 中断信号并发送给 \hyperref[mod:intmtrx]{中断矩阵}。

中断信号由 VDMA 的内部中断源生成。
\begin{itemize}
    \item 一般中断：
    \begin{itemize}
        \item \label{int:DMACSLVIFUNDEFINEDREGDECERRINT}DMAC\_SLVIF\_UNDEFINEDREG\_DEC\_ERR\_INT：从机接口未定义寄存器解码错误中断
        \item \label{int:DMACSLVIFCOMMONREGRD2WOERRINT}DMAC\_SLVIF\_COMMONREG\_RD2WO\_ERR\_INT：从机接口一般寄存器读入只写错误中断
        \item \label{int:DMACSLVIFCOMMONREGWR2ROERRINT}DMAC\_SLVIF\_COMMONREG\_WR2RO\_ERR\_INT：从机接口一般寄存器写到只读错误中断
        \item \label{int:DMACSLVIFCOMMONREGDECERRINT}DMAC\_SLVIF\_COMMONREG\_DEC\_ERR\_INT：从机接口一般寄存器解码错误中断
    \end{itemize}

    \item 通道中断 (\regindex{n} = 1 \verb+~+ 4)：

    \begin{itemize}
        \item \label{int:DMACCHNCHABORTEDINT}DMAC\_CH\regindex{n}\_CH\_ABORTED\_INT：通道中止中断
        \item \label{int:DMACCHNCHDISABLEDINT}DMAC\_CH\regindex{n}\_CH\_DISABLED\_INT：通道禁用中断
        \item \label{int:DMACCHNCHSUSPENDEDINT}DMAC\_CH\regindex{n}\_CH\_SUSPENDED\_INT：通道挂起中断
        \item \label{int:DMACCHNCHSRCSUSPENDEDINT}DMAC\_CH\regindex{n}\_CH\_SRC\_SUSPENDED\_INT：通道源挂起中断
        \item \label{int:DMACCHNSLVIFSHADOWREGWRONVALIDERRINT}DMAC\_CH\regindex{n}\_SLVIF\_SHADOWREG\_WRON\_VALID\_ERR\_INT：影子寄存器写有效错误中断
        \item \label{int:DMACCHNSLVIFWRONCHENERRINT}DMAC\_CH\regindex{n}\_SLVIF\_WRONCHEN\_ERR\_INT：从机接口写启用通道错误中断
        \item \label{int:DMACCHNSLVIFRD2RWOERRINT}DMAC\_CH\regindex{n}\_SLVIF\_RD2RWO\_ERR\_INT：从机接口读到只写错误中断
        \item \label{int:DMACCHNSLVIFWR2ROERRINT}DMAC\_CH\regindex{n}\_SLVIF\_WR2RO\_ERR\_INT：从机接口写入只读错误中断
        \item \label{int:DMACCHNSLVIFDECERRINT}DMAC\_CH\regindex{n}\_SLVIF\_DEC\_ERR\_INT：从机接口解码错误中断
        \item \label{int:DMACCHNSLVIFMULTIBLKTYPEERRINT}DMAC\_CH\regindex{n}\_SLVIF\_MULTIBLKTYPE\_ERR\_INT：从机接口多块类型错误中断
        \item \label{int:DMACCHNSHADOWREGORLLIINVALIDERRINT}DMAC\_CH\regindex{n}\_SHADOWREG\_OR\_LLI\_INVALID\_ERR\_INT：影子寄存器或 LLI 无效错误中断
        \item \label{int:DMACCHNLLIWRSLVERRINT}DMAC\_CH\regindex{n}\_LLI\_WR\_SLV\_ERR\_INT：LLI 写从机错误中断
        \item \label{int:DMACCHNLLIRDSLVERRINT}DMAC\_CH\regindex{n}\_LLI\_RD\_SLV\_ERR\_INT：LLI 读从机错误中断
        \item \label{int:DMACCHNLLIWRDECERRINT}DMAC\_CH\regindex{n}\_LLI\_WR\_DEC\_ERR\_INT：LLI 写解码错误中断
        \item \label{int:DMACCHNLLIRDDECERRINT}DMAC\_CH\regindex{n}\_LLI\_RD\_DEC\_ERR\_INT：LLI 读解码错误中断
        \item \label{int:DMACCHNDSTSLVERRINT}DMAC\_CH\regindex{n}\_DST\_SLV\_ERR\_INT：目标从机错误中断
        \item \label{int:DMACCHNSRCSLVERRINT}DMAC\_CH\regindex{n}\_SRC\_SLV\_ERR\_INT：源从机错误中断
        \item \label{int:DMACCHNDSTDECERRINT}DMAC\_CH\regindex{n}\_DST\_DEC\_ERR\_INT：目标解码错误中断
        \item \label{int:DMACCHNSRCDECERRINT}DMAC\_CH\regindex{n}\_SRC\_DEC\_ERR\_INT：源解码错误中断
        \item \label{int:DMACCHNDSTTRANSCOMPINT}DMAC\_CH\regindex{n}\_DST\_TRANSCOMP\_INT：目标传输完成中断
        \item \label{int:DMACCHNSRCTRANSCOMPINT}DMAC\_CH\regindex{n}\_SRC\_TRANSCOMP\_INT：源传输完成中断
        \item \label{int:DMACCHNDMATFRDONEINT}DMAC\_CH\regindex{n}\_DMA\_TFR\_DONE\_INT：DMA 传输完成中断
        \item \label{int:DMACCHNBLOCKTFRDONEINT}DMAC\_CH\regindex{n}\_BLOCK\_TFR\_DONE\_INT：块传输完成中断
    \end{itemize}
\end{itemize}
%%%%%%%%%%%%%%%%%%%%%%%%%%%%%%
%%% Programming Procedures %%%
%%%%%%%%%%%%%%%%%%%%%%%%%%%%%%

\section{配置流程}
\label{sec:vdma-prog-procedures}

\subsection{通用配置流程}

以下两个步骤是通用的配置流程，应在进行其他配置之前执行。

\begin{enumerate}
    \item 配置 \hyperref[fielddesc:HPSYSCLKRSTGDMASYSCLKEN]{HP\_SYS\_CLKRST\_GDMA\_SYS\_CLK\_EN} 为 1，使能 VDMA 的总线时钟。
    \item 通过 \hyperref[regdesc:TEEDMAGDMACHNWPMSREG]{TEE\_DMA\_GDMA\_CH\regindex{n}\_W\_PMS} 和 \hyperref[regdesc:TEEDMAGDMACHNRPMSREG]{TEE\_DMA\_GDMA\_CH\regindex{n}\_R\_PMS} 配置 VDMA 的访问权限。更多信息请查看章节 \ref{mod:permctrl} \textit{\nameref{mod:permctrl}}。
\end{enumerate}

\subsection{配置基于影子寄存器的多块传输}
\begin{enumerate}
    \item 读取 \hyperref[fielddesc:DMACCH1EN]{DMAC\_CH\regindex{n}\_EN} 选择一个可用的通道。

    \item 配置 \hyperref[fielddesc:DMACCHNSRCMULTBLKTYPE]{DMAC\_CH\regindex{n}\_SRC\_MULTBLK\_TYPE} 和/或 \hyperref[fielddesc:DMACCHNDSTMULTBLKTYPE]{DMAC\_CH\regindex{n}\_DST\_MULTBLK\_TYPE} 为 2，选择源传输和/或目标传输为基于影子寄存器的多块传输。

    \item 通过 \hyperref[regdesc:DMACCHNSAR0REG]{DMAC\_CH\regindex{n}\_SAR0\_REG} 和/或 \hyperref[regdesc:DMACCHNDAR0REG]{DMAC\_CH\regindex{n}\_DAR0\_REG}、\hyperref[fielddesc:DMACCHNBLOCKTS]{DMAC\_CH\regindex{n}\_BLOCK\_TS}、\hyperref[regdesc:DMACCHNCTL0REG]{DMAC\_CH\regindex{n}\_CTL0\_REG}、\hyperref[regdesc:DMACCHNCTL1REG]{DMAC\_CH\regindex{n}\_CTL1\_REG} 配置第一个块。VDMA 将这些值加载到相应的影子寄存器。

    \textbf{注意}：必须最后再将 \hyperref[regdesc:DMACCHNCTL1REG]{DMAC\_CH\regindex{n}\_CTL1\_REG} 的 \hyperref[fielddesc:DMACCHNSHADOWREGORLLIVALID]{DMAC\_CH\regindex{n}\_SHADOWREG\_OR\_LLI\_VALID} 位配置为 1。


    \item 配置 \hyperref[fielddesc:DMACCH1EN]{DMAC\_CH\regindex{n}\_EN} 为 1，使能相应通道。

    \item VDMA 根据配置发起 DMA 块传输操作。

        \begin{itemize}
        \item 块传输可能立即启动，也可能在握手请求之后启动，这取决于 \hyperref[fielddesc:DMACCHNTTFC]{DMAC\_CH\regindex{n}\_TT\_FC}。具体为，
        \begin{itemize}
            \item 当源是存储器的时候，源传输就立即启动；当目标是存储器的时候，目标传输就立即启动。
            \item 当源是外设时，从源获取数据需要等待与源的握手之后；当目标是外设时，往目标写入数据需要等待与目标的握手之后。
        \end{itemize}

        \item VDMA 检查 \hyperref[fielddesc:DMACCHNSHADOWREGORLLIVALID]{DMAC\_CH\regindex{n}\_SHADOWREG\_OR\_LLI\_VALID} 的值，如果为 0，VDMA 等待 \hyperref[fielddesc:DMACCHNBLKTFRRESUMEREQ]{DMAC\_CH\regindex{n}\_BLK\_TFR\_RESUMEREQ} 置 1 来指示影子寄存器可用，然后再次尝试影子寄存器获取操作。在这种情况下，VDMA 可能生成 \hyperref[int:DMACCHNSHADOWREGORLLIINVALIDERRINT]{DMAC\_CH\regindex{n}\_SHADOWREG\_OR\_LLI\_INVALID\_ERR\_INT} 中断。
        \item VDMA 检查 \hyperref[fielddesc:DMACCHNSHADOWREGORLLIVALID]{DMAC\_CH\regindex{n}\_SHADOWREG\_OR\_LLI\_VALID} 的值，如果为 1，VDMA 将影子寄存器内容复制到用于执行 DMA 传输的寄存器中（\hyperref[regdesc:DMACCHNSAR0REG]{DMAC\_CH\regindex{n}\_SAR0\_REG} 和/或 \hyperref[regdesc:DMACCHNDAR0REG]{DMAC\_CH\regindex{n}\_DAR0\_REG}、\hyperref[fielddesc:DMACCHNBLOCKTS]{DMAC\_CH\regindex{n}\_BLOCK\_TS}、\hyperref[regdesc:DMACCHNCTL0REG]{DMAC\_CH\regindex{n}\_CTL0\_REG} 和 \hyperref[regdesc:DMACCHNCTL1REG]{DMAC\_CH\regindex{n}\_CTL1\_REG}），并且清零 \hyperref[fielddesc:DMACCHNSHADOWREGORLLIVALID]{DMAC\_CH\regindex{n}\_SHADOWREG\_OR\_LLI\_VALID}。
            \begin{itemize}
            \item 如果 VDMA 读到 \hyperref[fielddesc:DMACCHNSHADOWREGORLLILAST]{DMAC\_CH\regindex{n}\_SHADOWREG\_OR\_LLI\_LAST} 为 1，则认为当前块是最后一块，并在当前块传输结束时完成 DMA 传输操作。
            \item 如果 VDMA 读到 \hyperref[fielddesc:DMACCHNSHADOWREGORLLILAST]{DMAC\_CH\regindex{n}\_SHADOWREG\_OR\_LLI\_LAST} 为 0，则认为还有一个或多个块需要传输，并在当前块传输结束时再次检查 \hyperref[fielddesc:DMACCHNSHADOWREGORLLIVALID]{DMAC\_CH\regindex{n}\_SHADOWREG\_OR\_LLI\_VALID}。
        \end{itemize}
    \end{itemize}

    \item 软件等待 \hyperref[fielddesc:DMACCHNSHADOWREGORLLIVALID]{DMAC\_CH\regindex{n}\_SHADOWREG\_OR\_LLI\_VALID} 自动清零。

    \item 通过 \hyperref[regdesc:DMACCHNSAR0REG]{DMAC\_CH\regindex{n}\_SAR0\_REG} 和/或 \hyperref[regdesc:DMACCHNDAR0REG]{DMAC\_CH\regindex{n}\_DAR0\_REG}、\hyperref[fielddesc:DMACCHNBLOCKTS]{DMAC\_CH\regindex{n}\_BLOCK\_TS}、\hyperref[regdesc:DMACCHNCTL0REG]{DMAC\_CH\regindex{n}\_CTL0\_REG}、\hyperref[regdesc:DMACCHNCTL1REG]{DMAC\_CH\regindex{n}\_CTL1\_REG} 配置下一个块的信息。若下一个块是最后一个块，需要将 \hyperref[fielddesc:DMACCHNSHADOWREGORLLILAST]{DMAC\_CH\regindex{n}\_SHADOWREG\_OR\_LLI\_LAST} 配置为 1。

    \textbf{注意}：必须最后再将 \hyperref[fielddesc:DMACCHNSHADOWREGORLLIVALID]{DMAC\_CH\regindex{n}\_SHADOWREG\_OR\_LLI\_VALID} 配置为 1。

    \item 软件等待 \hyperref[int:DMACCHNBLOCKTFRDONEINT]{DMAC\_CH\regindex{n}\_BLOCK\_TFR\_DONE\_INT} 中断，或者轮询块传输完成状态指示位 \hyperref[fielddesc:DMACCHNBLOCKTFRDONEINTSTAT]{DMAC\_CH\regindex{n}\_BLOCK\_TFR\_DONE\_INTSTAT} 直到读到 1。之后回到步骤 6 循环。


\end{enumerate}

\textbf{说明：}
如果产生了 \hyperref[int:DMACCHNSHADOWREGORLLIINVALIDERRINT]{DMAC\_CH\regindex{n}\_SHADOWREG\_OR\_LLI\_INVALID\_ERR\_INT} 中断，推荐按照以下步骤恢复传输：

\begin{itemize}
    \item 通过 \hyperref[regdesc:DMACCHNSAR0REG]{DMAC\_CH\regindex{n}\_SAR0\_REG} 和/或 \hyperref[regdesc:DMACCHNDAR0REG]{DMAC\_CH\regindex{n}\_DAR0\_REG}、\hyperref[fielddesc:DMACCHNBLOCKTS]{DMAC\_CH\regindex{n}\_BLOCK\_TS}、\hyperref[regdesc:DMACCHNCTL0REG]{DMAC\_CH\regindex{n}\_CTL0\_REG} 和 \hyperref[regdesc:DMACCHNCTL1REG]{DMAC\_CH\regindex{n}\_CTL1\_REG} 配置下一个块。
    \item 配置 \hyperref[fielddesc:DMACCHNCLEARSHADOWREGORLLIINVALIDERRINTSTAT] {DMAC\_CH\regindex{n}\_CLEAR\_SHADOWREG\_OR\_LLI\_INVALID\_ERR\_INTSTAT} 清除中断。
    \item 配置 \hyperref[fielddesc:DMACCHNBLKTFRRESUMEREQ]{DMAC\_CH\regindex{n}\_BLK\_TFR\_RESUMEREQ} 为 1 请求恢复块传输。
\end{itemize}

\subsection{配置基于链表的多块传输}
\begin{enumerate}
    \item 读取 \hyperref[fielddesc:DMACCH1EN]{DMAC\_CH\regindex{n}\_EN} 选择一个可用的通道。
    \item 配置 \hyperref[fielddesc:DMACCHNSRCMULTBLKTYPE]{DMAC\_CH\regindex{n}\_SRC\_MULTBLK\_TYPE} 和/或 \hyperref[fielddesc:DMACCHNDSTMULTBLKTYPE]{DMAC\_CH\regindex{n}\_DST\_MULTBLK\_TYPE} 为 3，选择源传输和/或目标传输为基于链表的多块传输。
    \item 通过 \hyperref[regdesc:DMACCHNLLP0REG]{DMAC\_CH\regindex{n}\_LLP0\_REG} 配置第一个链表项的基地址，以及链表所在的主机接口。
    \item 在系统存储器中创建一个或多个链表项。可以提前创建这个链表项或者使用 \hyperref[fielddesc:DMACCHNSHADOWREGORLLIVALID]{DMAC\_CH\regindex{n}\_SHADOWREG\_OR\_LLI\_VALID} 和 \hyperref[fielddesc:DMACCHNSHADOWREGORLLILAST]{DMAC\_CH\regindex{n}\_SHADOWREG\_OR\_LLI\_LAST} 动态地扩展链表。动态扩展链表的具体操作为：
    \begin{itemize}
        \item 在 DMA 传输过程中，当链表没有准备好时，软件可以把链表中 \hyperref[fielddesc:DMACCHNSHADOWREGORLLIVALID]{DMAC\_CH\regindex{n}\_SHADOWREG\_OR\_LLI\_VALID} 置 0，然后等待准备好链表后，再置 1，用于拓展链表。
        \item 当要结束 DMA 传输时，把最后一个链表的 \hyperref[fielddesc:DMACCHNSHADOWREGORLLILAST]{DMAC\_CH\regindex{n}\_SHADOWREG\_OR\_LLI\_LAST} 置 1，来表示传输结束。
    \end{itemize}
    \item 配置 \hyperref[fielddesc:DMACCH1EN]{DMAC\_CH\regindex{n}\_EN} 为 1，使能相应通道。

\begin{tiplisting}
步骤 4 和 5 的顺序可以互换。但是，如果先执行步骤 5，或者在多块传输期间，下一个块传输的链表项在系统存储器中不可用（即获取的链表项的 \hyperref[fielddesc:DMACCHNSHADOWREGORLLIVALID]{DMAC\_CH\regindex{n}\_SHADOWREG\_OR\_LLI\_VALID} 为 0），则 VDMA 可能会产生 \hyperref[int:DMACCHNSHADOWREGORLLIINVALIDERRINT]{DMAC\_CH\regindex{n}\_SHADOWREG\_OR\_LLI\_INVALID\_ERR\_INT} 中断。
\end{tiplisting}

    \item VDMA 根据配置发起 DMA 块传输。块传输可能立即启动，也可能在硬件握手请求之后启动，这取决于 \hyperref[fielddesc:DMACCHNTTFC]{DMAC\_CH\regindex{n}\_TT\_FC}。VDMA 将链表内容复制到用于执行 DMA 传输的寄存器中（\hyperref[regdesc:DMACCHNSAR0REG]{DMAC\_CH\regindex{n}\_SAR0\_REG} 和/或 \hyperref[regdesc:DMACCHNDAR0REG]{DMAC\_CH\regindex{n}\_DAR0\_REG}、\hyperref[fielddesc:DMACCHNBLOCKTS]{DMAC\_CH\regindex{n}\_BLOCK\_TS}、\hyperref[regdesc:DMACCHNCTL0REG]{DMAC\_CH\regindex{n}\_CTL0\_REG} 和 \hyperref[regdesc:DMACCHNCTL1REG]{DMAC\_CH\regindex{n}\_CTL1\_REG}），然后发起 DMA 块传输。

    \item 在链表获取阶段：

        \begin{itemize}
        \item 如果 VDMA 读到 \hyperref[fielddesc:DMACCHNSHADOWREGORLLILAST]{DMAC\_CH\regindex{n}\_SHADOWREG\_OR\_LLI\_LAST} 为 1，则认为当前块是最后一个块，并在当前块传输结束时完成 DMA 传输。
        \item 如果 VDMA 读到 \hyperref[fielddesc:DMACCHNSHADOWREGORLLILAST]{DMAC\_CH\regindex{n}\_SHADOWREG\_OR\_LLI\_LAST} 为 0，则认为还有一个或多个块需要传输，并且回到步骤 6。
        \item 如果 VDMA 读到 \hyperref[fielddesc:DMACCHNSHADOWREGORLLIVALID]{DMAC\_CH\regindex{n}\_SHADOWREG\_OR\_LLI\_VALID} 为 0，则可能产生 \hyperref[int:DMACCHNSHADOWREGORLLIINVALIDERRINT]{DMAC\_CH\regindex{n}\_SHADOWREG\_OR\_LLI\_INVALID\_ERR\_INT} 中断。VDMA 等待直到软件将 \hyperref[fielddesc:DMACCHNBLKTFRRESUMEREQ]{DMAC\_CH\regindex{n}\_BLK\_TFR\_RESUMEREQ} 置 1 来指示链表的可用性，然后再尝试读取链表。
    \end{itemize}
\end{enumerate}

\subsection{配置单块传输}

\begin{enumerate}
    \item 读取 \hyperref[fielddesc:DMACCH1EN]{DMAC\_CH\regindex{n}\_EN} 选择一个可用的通道。
    \item 配置 \hyperref[fielddesc:DMACCHNSRCMULTBLKTYPE]{DMAC\_CH\regindex{n}\_SRC\_MULTBLK\_TYPE} 和 \hyperref[fielddesc:DMACCHNDSTMULTBLKTYPE]{DMAC\_CH\regindex{n}\_DST\_MULTBLK\_TYPE} 为 0，选择连续地址的单块传输。
    \item 配置 \hyperref[regdesc:DMACCHNSAR0REG]{DMAC\_CH\regindex{n}\_SAR0\_REG} 和/或 \hyperref[regdesc:DMACCHNDAR0REG]{DMAC\_CH\regindex{n}\_DAR0\_REG}、 \hyperref[fielddesc:DMACCHNBLOCKTS]{DMAC\_CH\regindex{n}\_BLOCK\_TS}、\hyperref[regdesc:DMACCHNCTL0REG]{DMAC\_CH\regindex{n}\_CTL0\_REG} 和 \hyperref[regdesc:DMACCHNCTL1REG]{DMAC\_CH\regindex{n}\_CTL1\_REG}。
    \item 配置 \hyperref[fielddesc:DMACCH1EN]{DMAC\_CH\regindex{n}\_EN} 为 1 使能相应通道。

    \item 软件等待 \hyperref[int:DMACCHNBLOCKTFRDONEINT]{DMAC\_CH\regindex{n}\_BLOCK\_TFR\_DONE\_INT} 中断，或轮询块传输完成状态指示位 \hyperref[fielddesc:DMACCHNBLOCKTFRDONEINTSTAT]{DMAC\_CH\regindex{n}\_BLOCK\_TFR\_DONE\_INTSTAT} 直到读到 1。
\end{enumerate}

% >>> If your module has no registers, delete sample text
%       for Sections Register Summary and Registers below <<<

%%%%%%%%%%%%%%%%%%%%%%%%
%%% Register Summary %%%
%%%%%%%%%%%%%%%%%%%%%%%%

% >>> Update the sample text below and insert
%     the info on register summary into the subfile <<<

\newpage
\hypertarget{VDMA-reg-summ}{}
\section{寄存器列表}
\label{sec:VDMA-reg-summ}

本小节的所有地址均为相对于 VDMA 控制器基地址的地址偏移量（相对地址），具体基地址请见章节 \ref{mod:sysmem} \textit{\nameref{mod:sysmem}} 中的表 \ref{tab:sysmem-base-address}。

请查看章节 \hyperref[glossary-access-types]{\textit{寄存器的访问类型}}，了解“访问”列缩写的含义。

\subfile{\modulefiles/06-VDMA-reg-summ__CN}


%%%%%%%%%%%%%%%%%
%%% Registers %%%
%%%%%%%%%%%%%%%%%

% >>> Update the sample text below and insert
%      the info on registers into the subfile <<<

\newpage
\section{寄存器}

本小节的所有地址均为相对于 VDMA 控制器基地址的地址偏移量（相对地址），具体基地址请见章节 \ref{mod:sysmem} \textit{\nameref{mod:sysmem}} 中的表 \ref{tab:sysmem-base-address}。

关于保留 (reserved) 域的处理，请查看章节 \hyperref[programming-reserved-field]{如何配置寄存器的保留域}。

\subfile{\modulefiles/06-VDMA-reg__CN}


\end{document}
